\section{Miscellanious}

\subsection{Logic}
We let $P$ and $Q$ be propositions.

\begin{itemize}
\item $P\Rightarrow Q \iff \neg P\text{ or }Q$
\item $\neg (P\Rightarrow Q) \iff P\text{ and } \neg Q$
\item Contrapositive: $P\Rightarrow Q \iff \neg Q\Rightarrow \neg P$
\end{itemize}

\subsection{Fundamental Groups}

Some random fundamental groups to remember. Note that fundamental groups 
are usually taken with respect to a base point, but I'm lazy so I haven't written this below.

\begin{proposition}
If $X$ is path-connected, then $\pi_1(X, x_0)$ is independent of the 
base point.
\end{proposition}

\begin{proposition}
$\pi_1(X \times Y, (x_0, y_0)) \cong \pi_1(X, x_0) \times \pi_1(Y, y_0)$.
\end{proposition}

\begin{example}
$\pi_1(\T^2) \cong \Z^2$, because $\T^2 = \S^1 \times \S^1$.
\end{example}

\begin{remark}
If $X, Y$ are path connected, we can just omit base point.
\end{remark}

\begin{df}
A topological space $X$ is \emph{simply connected} if $\pi_1(X) = 0$, or $\{e\}$.
\end{df}

\begin{itemize}
    \item $\pi_1(\mathbb{R}^n)=\{e\}$
    \item $\pi_1((0,1))=\{e\}$
    \item $\pi_1([0,1])=\{e\}$
    \item $\pi_1(D^n)=\{e\}$
    \item $\pi_1(\S^1)=\mathbb{Z}$
    \item $\pi_1(\S^n)=\{e\}$ for $n\geq2$
    \item $\pi_1(\mathbb{R}^2\setminus\{0\})=\mathbb{Z}$
    \item $\pi_1(\mathbb{R}^n\setminus\{0\})=\{e\}$ for $n\geq3$
    \item $\pi_1(S^1\times S^1)=\mathbb{Z}^2$
    \item $\pi_1(\mathbb{R}^2\setminus\{p_1,\ldots,p_k\})=F_k$
    \item $\pi_1(\mathbb{RP}^n)=\mathbb{Z}/2\mathbb{Z}$ for $n\geq2$
    \item $\pi_1(\text{Klein bottle})
    =\langle a,b\mid aba^{-1}=b^{-1}\rangle$
    \item $\pi_1\left(\bigvee_{i=1}^k S^1\right)=F_k$, the free group on $k$ generators.
    \item $\pi_1(X \times Y, ) \cong \pi_1(X) \times \pi_1(Y)$.
\end{itemize}

\subsection{Some Specific Rings}

\begin{proposition}
Let $R$ be an idempotent commutative ring, i.e. $a^2 = a$, for all $a \in R$.
Then, the following hold:
\begin{itemize}
    \item $\char(R) = 2$.
    \item $R$ is commutative.
    \item If $R$ is an integral domain, then $R \cong \F_2$.
\end{itemize}
\end{proposition}

\begin{example}
Let $R$ be Noetherian. $R[x_1, x_2, \dots]$ is not Noetherian (as a ring).
\end{example}

\begin{remark}
Note that if $R$ is a ring with only two left and/ or right ideals, 
then $R$ is not necessarily a field. $R$ is merely a division ring, but 
not commutative. Recall that fields are \emph{commutative division rings}.
\end{remark}

\begin{example}[non-Commutative division ring]
Hurwitz quaternions are a great exmaple of a noncommutative division ring.
This is $\H = \R \oplus \R i \oplus \R j \oplus \R k$ subject to the usual relations.
More generally, quaternion algebras over fields give good examples.
\end{example}

\begin{remark}
An interesting result of Frobenius states that every f.d. associative 
real division algebra is isomorphic to $\R$, $\C$, or $\H$.
\end{remark}